% Responsable: Jerónimo
% Fuentes: rama claude/compat-10nov, out10/ (corridas del 10-nov, escenario común, 30 réplicas).
% Las tablas se generan desde los CSV (tablas/*.tex); no editar los números a mano.
\subsection{Información imperfecta (Modelos 3 y 4)}\label{sec:res34}

Todas las cifras de esta sección son del escenario común
(Sección~\ref{sec:diseno}). El criterio de costo usa \cReplicasCosto{} réplicas
del sorteo por meta, y las mismas réplicas para M3, M4 y el placebo. De las
\cCorridasCosto{} corridas del criterio de costo, todas terminaron con
optimalidad demostrada y brecha menor que $10^{-6}$.

\subsubsection{Costo de cumplir la meta}

La Tabla~\ref{tab:costos} resume el costo de reducir $\Mcero$ en metas del 2 al
30\,\%. Se lee de izquierda a derecha como la escalera de información:
\begin{itemize}
\item el costo con información perfecta (Modelo 2);
\item lo que cuesta la misma meta cuando el Estado solo conoce las zonas (M3) o
  las zonas y el tipo de hogar (M4);
\item el placebo, que tiene la misma partición que el M4 pero sin contenido
  demográfico.
\end{itemize}

\begin{table}[ht]
\centering\footnotesize
\caption{Costo de cumplir la meta de reducción de $\Mcero$. IP: información perfecta.
Montos en miles de millones de \$ de 2025, mediana de \cReplicasCosto{} réplicas. M4/placebo es la mediana del
cociente pareado réplica a réplica, con su IC del 95\,\% por \emph{bootstrap}.
Inclusión: fracción del gasto que llega a hogares no pobres (M3 / M4).}
\label{tab:costos}
\begin{tabular}{@{}rrrrrrrcc@{}}
\toprule
& \multicolumn{4}{c}{Costo} & \multicolumn{2}{c}{Veces la IP} & & \\
\cmidrule(lr){2-5}\cmidrule(lr){6-7}
Meta & IP (M2) & M3 & M4 & Placebo & M3 & M4 & M4/placebo & Inclusión \\
\midrule
\csname @@input\endcsname tablas/costos.tex 
\bottomrule
\end{tabular}
\end{table}

Tres resultados se destacan.

\paragraph{Focalizar sin información individual cuesta varias veces más.} Con
metas modestas, del 2 al 10\,\%, el M3 cuesta entre \cMtresIPmin{} y
\cMtresIPmax{} veces lo que costaría con información perfecta. Con metas del
20 al 30\,\% la razón baja a unas \cMtresIPalta{} veces. La brecha se cierra con
metas grandes porque para cumplirlas hay que llegar a casi todos los pobres
de todas formas: la información sirve para elegir a quién atender primero, y
pierde valor cuando hay que atender a casi todos.

\paragraph{El M4 ahorra, pero hasta metas del 18\,\% el ahorro viene de la
partición.} El M4 es más barato que el M3 en todas las metas salvo la del 16\,\%, y el
intervalo del cociente pareado M4/M3 queda por debajo de 1 en 13 de las 15 metas. Sin embargo,
hasta la meta del 18\,\% su costo es en general indistinguible del placebo: el
cociente pareado M4/placebo ronda 1 y su intervalo contiene el 1 en 6 de las 9
metas, y en las otras tres la diferencia es pequeña y va en ambos sentidos (0,90
con la meta del 6\,\% y 1,04 con las del 16 y 18\,\%). En ese rango la ventaja
del M4 sobre el M3 se debe a repartir los indicadores privados en más celdas, no
a la información demográfica.

La Tabla~\ref{tab:gasto} muestra por qué. Hasta la meta del 18\,\% los tres
modelos gastan entre el 97 y el 100\,\% del presupuesto en salud (aseguramiento
y barreras de acceso). Esas privaciones varían entre tipos de hogar tan poco como
entre zonas (Tabla~\ref{tab:discrimina}), así que focalizarlas por tipo no rinde
más que una partición al azar. La educación entra desde la meta del
20\,\% y con la del 30\,\% se lleva entre el 68 y el 81\,\% del gasto. Ahí la
demografía sí focaliza, porque las privaciones educativas se concentran en
ciertos tipos de hogar: la inasistencia escolar, por ejemplo, es
\cDiscInasClu{} veces más frecuente en el tipo más privado que en el menos
privado. Desde metas del 22\,\% el M4 cuesta entre un \cMcuatroPlacMin{} y un
\cMcuatroPlacMax{} menos que el placebo. En ningún modelo el gasto en niñez y
vivienda pasa del 2\,\%, y el de trabajo no pasa del 13\,\%.

\begin{table}[ht]
\centering\footnotesize
\caption{En qué se gasta: participación promedio de la salud, la educación y la
zona de cabecera en el gasto (\%, \cReplicasCosto{} réplicas).}
\label{tab:gasto}
\begin{tabular}{@{}rccccccccc@{}}
\toprule
& \multicolumn{3}{c}{Salud} & \multicolumn{3}{c}{Educación} & \multicolumn{3}{c}{Cabecera} \\
\cmidrule(lr){2-4}\cmidrule(lr){5-7}\cmidrule(lr){8-10}
Meta & M3 & M4 & Plac. & M3 & M4 & Plac. & M3 & M4 & Plac. \\
\midrule
\csname @@input\endcsname tablas/gasto.tex
\bottomrule
\end{tabular}
\end{table}

\paragraph{El M3 se concentra en la cabecera.} Con metas del 6 al 14\,\%, el M3
pone entre el 84 y el 97\,\% del gasto en la cabecera, frente a entre el 58 y el
73\,\% del M4 y el placebo (Tabla~\ref{tab:gasto}). En la cabecera vive más de
la mitad de las personas pobres del departamento, pero la incidencia es la más
baja (\cHcabecera{}). Al comprar cobertura en esa zona, la mayor parte de los
privados atendidos no es pobre, y eso coincide con que el error de inclusión del
M3 sea más alto que el del M4 justo en esas metas.

\paragraph{Los errores de focalización son grandes.} Entre el 25 y el 69\,\% del
gasto llega a hogares que no eran pobres. La inclusión es máxima con metas
intermedias, cuando el modelo compra cobertura en zonas donde la mayoría de los
privados no es pobre. Al mismo tiempo, con metas del 2\,\% el 97\,\% de las
personas pobres no recibe ninguna intervención, y con metas del 30\,\%, el 58\,\%.

\subsubsection{Personas que salen de la pobreza con un presupuesto dado}

Con el criterio de equidad (Tabla~\ref{tab:equidad}), el M3 saca de la pobreza
entre el \cMtresPctMin{} y el \cMtresPctMax{} de las personas que saldrían con
información perfecta, y el M4 entre el 19 y el 67\,\%. Con $B=\$2$~mm la
diferencia entre M4 y M3 no es significativa. Con $B=\$5$~mm el M4 saca 1,86
veces las personas del M3, aunque el placebo logra lo mismo. Con $B\ge\$50$~mm el
M4 supera al placebo en un \cMcuatroPlacEq{}, con intervalos que excluyen el 1.
Las restricciones de severidad del Modelo 1 ($\rho=0{,}10$) son infactibles con
$B\le\$10$~mm en las \cInfactiblesRho{} corridas de M3, M4 y el placebo, y
factibles en todas con $B\ge\$25$~mm.

\subsubsection{Sensibilidad a $k$}

La elección de $k$ se apoyó en un análisis con 50 réplicas pareadas para
$k\in\{2,3,4,5,6,8\}$. Ese análisis se hizo con una versión simplificada de los
costos (una intervención por hogar, $u_{ij}=1$), así que lo usamos solo por su
forma, no por sus cifras. Muestra tres cosas. Primero, con $k=2$ el M4 saca de la
pobreza \emph{menos} personas que el M3 (0,71 veces con $B=\$10$~mm): la
partición separa los hogares por tamaño, que casi no se relaciona con la pobreza.
Segundo, de $k=3$ a $k=4$ la ganancia es grande, y desde ahí se aplana (1,92 veces
el M3 con $k=4$ y 1,93 con $k=5$). Tercero, con $k=8$ el M4 mejora otra vez, pero
lo hace con 26 celdas de menos de 10 hogares privados: está simulando información
individual más que focalización por tipo. La lección de política es que la
información demográfica solo ayuda si está correlacionada con las privaciones
que se quieren atacar.
